\section{Why population stability alone does not complete the argument}
\label{sec:distributional-only-obstruction}

Small population instability for each network does not itself control
instability at training inputs selected by posterior conditioning. This
is a statistical obstruction, even when compression is trivial away from
the exceptional inputs. The example below concerns an abstract prior;
it is not a lower bound for the particular Gaussian network prior.

\begin{proposition}[Sparse exceptional sets can absorb the training sample]
\label{prop:distributional-sparse-patches}
Let $\mathcal X$ have $M$ points with the uniform distribution, let
$f\equiv+1$, and fix $1\le r\le M/2$. Write $\rho=r/M$.
For every subset $B\subset\mathcal X$ of size $r$, let $h_B$ be $+1$
on $B$ and $-1$ elsewhere, and assign it exceptional set $B$.
The target has empty exceptional set. Put prior mass $q\in(0,1)$ on
$f$, and distribute mass $1-q$ uniformly over the $h_B$.
Then every hypothesis has population exception probability at most
$\rho$ and has an exact compression outside its exceptions into the
two-element class $\{+1,-1\}$.

For every labeled sample of size $n\le r/2$, if
\[
 q\le\tfrac12(\rho/2)^n,
\]
the ordinary consistency posterior has Gibbs error at least
$(1-\rho)/2$. The assertion holds for every sample, including samples
with repeated inputs.
\end{proposition}

\begin{proof}
Let $s\le n$ be the number of distinct observed inputs. A bad hypothesis
$h_B$ is consistent exactly when all these inputs belong to $B$.
Under the uniform prior on size-$r$ subsets, this has probability
\[
 p_s=\frac{\binom{M-s}{r-s}}{\binom Mr}
     =\prod_{j=0}^{s-1}\frac{r-j}{M-j}
     \ge(\rho/2)^s\ge(\rho/2)^n,
\]
because $s\le r/2$. Therefore the posterior mass of bad hypotheses is
\[
 \frac{(1-q)p_s}{q+(1-q)p_s}\ge\frac12.
\]
Indeed $q\le p_s/2$ and $q\le1/2$ imply $(1-q)p_s\ge q$.
Every bad hypothesis has true error $1-\rho$. It also has empirical
exception probability exactly one, although its population exception
probability is only $\rho$. The two-element compressor simply maps
$h_B$ to the constant negative function and leaves the target positive.
\end{proof}

If $q=e^{-aD}$, where $a=N/t^2$ is the width-and-scale factor in the
Gaussian target-mass bound, the displayed condition follows whenever
\[
 aD\ge n\log(2/\rho)+\log2.
\]
Thus a target-mass estimate of this form, a width-independent finite
compression class, and per-hypothesis population instability at most
$\rho$ do not by themselves imply a width-independent sample bound.
An additional structural property of the actual Gaussian prior could
still do so; the proposition neither supplies nor rules out such a
property.

This obstruction is relevant to interpreting certificates. An independent
held-out set can estimate the population instability of one fixed trained
network, or of one frozen posterior. That estimate does not bound the
training instability of a posterior obtained by selecting among many
possible exceptional sets. The controlled-family localized theorem
removes exactly this ambiguity by paying for the family of exceptional
sets. A trusted input region fixed before labeled training is the
simplest special case.

One can also require stability at every observed training input in an
input-dependent prior while retaining the ordinary label likelihood.
That gives the same mathematical benefit as a robust training constraint:
the empirical exception probability is then zero. It changes the posterior
conditioning event, however, and should not be presented as a consequence
of population stability alone under the original event.
